\begin{helper}
Let \(N\) be a natural number, let
\(v\in\mathbb R^{\operatorname{Fin}(N)}\), and let \(L\in\mathbb R\).
Define \(h\colon\mathbb R^{\operatorname{Fin}(N)}\to\mathbb R\) by
\[
h(y)=\frac{1}{N!}\sum_{\pi\in S_N}
\exp\!\left(L\sum_{i=1}^N y_i v_{\pi(i)}\right).
\]
Then \(h\) is convex in the following pointwise sense:
for all \(y,z\) and all \(0\le\theta\le1\),
\[
h((1-\theta)y+\theta z)\le (1-\theta)h(y)+\theta h(z).
\]
It is also symmetric: for every permutation \(\sigma\) of
\(\operatorname{Fin}(N)\),
\[
h(y\circ\sigma)=h(y).
\]
\end{helper}
\begin{proof}
Fix a permutation \(\pi\in S_N\), and define
\[
\ell_\pi(y)=L\sum_i y_i v_{\pi(i)}.
\]
This is an affine, indeed linear, function of the vector \(y\).  For
\(0\le\theta\le1\),
\[
\ell_\pi((1-\theta)y+\theta z)
=(1-\theta)\ell_\pi(y)+\theta\ell_\pi(z),
\]
by distributing the finite sum and the scalar multiplications.  The
exponential function on \(\mathbb R\) is convex, so
\[
\exp(\ell_\pi((1-\theta)y+\theta z))
\le
(1-\theta)\exp(\ell_\pi(y))+\theta\exp(\ell_\pi(z)).
\]
Summing this inequality over all \(\pi\in S_N\) and multiplying by the
nonnegative scalar \((N!)^{-1}\) gives
\[
h((1-\theta)y+\theta z)
\le (1-\theta)h(y)+\theta h(z),
\]
because finite sums distribute over addition and scalar multiplication.
This proves convexity.

For symmetry, fix a permutation \(\sigma\).  Then
\[
h(y\circ\sigma)
=\frac1{N!}\sum_{\pi\in S_N}
\exp\!\left(L\sum_i y_{\sigma(i)}v_{\pi(i)}\right).
\]
In the inner sum, put \(j=\sigma(i)\), equivalently
\(i=\sigma^{-1}(j)\).  Since \(\sigma\) is a bijection,
\[
\sum_i y_{\sigma(i)}v_{\pi(i)}
=\sum_j y_j\,v_{\pi(\sigma^{-1}(j))}.
\]
As \(\pi\) ranges over all permutations, the composite
\(\pi\circ\sigma^{-1}\) also ranges over all permutations exactly once.
Thus reindexing the finite outer sum by
\(\tau=\pi\circ\sigma^{-1}\) gives
\[
h(y\circ\sigma)
=\frac1{N!}\sum_{\tau\in S_N}
\exp\!\left(L\sum_j y_j v_{\tau(j)}\right)
=h(y).
\]
\end{proof}
